\section{Introduction}
\label{sec:intro}

Apache Spark is a widely used engine for batch data pipelines in cloud data platforms, and
many such pipelines are \emph{recurring}: the same job runs every hour, day or week over
fresh data. Spark exposes many optimization mechanisms for such pipelines, from execution
parameters such as the number of shuffle partitions to storage-level maintenance such as
compacting many small files into fewer large ones. A substantial body of research automates
these mechanisms, either by searching the configuration space of the engine
\cite{nguyen2018tuning,xin2022locat,lyu2024optimizer,herodotou2020survey} or by deciding how
tables should be laid out and maintained \cite{gruenheid2025autocomp,meduri2026pto}.

Optimization, however, is not free. Searching for a better shuffle-partition setting requires
executing the pipeline under candidate settings, and compaction is itself a Spark job that
reads and rewrites the data. An action that makes one execution faster is therefore not
automatically economical: its one-time cost must be recovered by the savings of the executions
that follow. Whether that happens depends on the size of the saving, on the cost of the action,
and on how many times the pipeline will run again. A pipeline that runs once should rarely be
tuned; a pipeline that runs thousands of times almost always should. Between those extremes,
the decision is often made without an explicit cost comparison, for example by leaving Spark's
defaults in place or by applying an optimization unconditionally.

Prior work has established each ingredient of a better answer in isolation. Spark tuners treat
their own overhead as a search-efficiency cost to be minimized rather than as a term to be
recovered \cite{xin2022locat,li2025loftune}. Workload characteristics are an accepted signal for
optimization \cite{shen2023rover,parayil2025workload}. Recent lakehouse systems weigh the benefit
of compaction or layout changes against their cost at production scale
\cite{gruenheid2025autocomp,meduri2026pto,nakandala2026autoliquid}, and Fekry
\emph{et al.}~\cite{fekry2020tune} showed that, for general configuration tuning, the tuning
cost must be amortized over future executions. What is missing is a controlled, empirical
evaluation that applies this break-even principle \emph{jointly} to two different categories
of Spark optimization (an execution-level action and a storage-level action) under one
explicit, interpretable decision rule, and that identifies the workload conditions and
execution frequencies at which each action becomes economically justified.

This paper addresses that gap with CAO-DF, a Cost-Aware Optimization Decision Framework. Given
a workload's characteristics and the expected number of future executions $N$, CAO-DF chooses
among three outcomes: keep the static configuration, tune the number of shuffle partitions
(action A1), or compact the input files (action A2). It does so by comparing total costs of
the form $\Copt + N \cdot \Crun$ and reporting the optimization break-even point (OBEP), the
smallest $N$ at which an action pays for itself. We do not claim that Spark tuning,
workload-aware optimization, adaptive layout or cloud cost optimization are new; our
contribution is the focused empirical evaluation of when two existing actions pay for
themselves. We study three research questions:
\begin{itemize}
  \item \textbf{RQ1:} How do data volume, file fragmentation, shuffle volume and execution
        frequency affect the cost and performance benefits of shuffle-partition tuning and
        file compaction?
  \item \textbf{RQ2:} Can a workload-aware cost-benefit framework reduce total workload cost
        compared with a static Spark configuration?
  \item \textbf{RQ3:} Under what workload conditions does applying shuffle-partition tuning or
        file compaction become economically justified?
\end{itemize}

\noindent The contributions of this paper are:
\begin{enumerate}
  \item a simple, interpretable cost-benefit decision framework for two Spark optimization
        actions of different kinds, with an explicit break-even criterion;
  \item an empirical characterization, from 250 measured executions under a randomized-block
        protocol, of how data volume, file fragmentation and shuffle volume affect the value
        of each action;
  \item a quantitative break-even analysis linking measured optimization overhead to future
        savings, with a sensitivity analysis over price and cost proxy; and
  \item practical guidance on when to apply an optimization and when to keep the static
        pipeline, together with a public, reproducible artifact.
\end{enumerate}

\tbd{One or two sentences stating the headline results with numbers, after Phases 8--10.}

The rest of the paper is organized as follows. Section~\ref{sec:related} reviews related work.
Section~\ref{sec:problem} defines the problem and the cost model, and
Section~\ref{sec:framework} presents the decision framework. Section~\ref{sec:method}
describes the experimental methodology, Section~\ref{sec:results} reports the results, and
Section~\ref{sec:discussion} discusses them. Section~\ref{sec:threats} states the threats to
validity, and Section~\ref{sec:conclusion} concludes.
